% OWNER: Member 1 (with input from Members 2 and 3). SKELETON — fill after final runs.
\section{Final comparison and ablations}\label{sec:comparison}

\begin{table}[h]
\centering\small
\caption{Final agents on development and held-out games (mean $\pm$ s.d., 3 seeds,
\todo{5000}-action budget). Official RHAE from the toolkit scorecard; Kaggle score
is the public leaderboard on hidden games.}\label{tab:final}
\begin{tabular}{lccccc}
\toprule
Agent & RHAE (dev) & RHAE (held-out) & Levels (held-out) & Actions to 1st level & Wall-clock \\
\midrule
Random            & \todo{} & \todo{} & \todo{} & \todo{} & \todo{} \\
DQN baseline      & \todo{} & \todo{} & \todo{} & \todo{} & \todo{} \\
DQN + imp1 + imp2 & \todo{} & \todo{} & \todo{} & \todo{} & \todo{} \\
PPO baseline      & \todo{} & \todo{} & \todo{} & \todo{} & \todo{} \\
PPO + imp1 + imp2 & \todo{} & \todo{} & \todo{} & \todo{} & \todo{} \\
\bottomrule
\end{tabular}
\end{table}

\paragraph{Ablations.} \todo{For each algorithm: full agent, minus imp1, minus
imp2. Which component accounts for the gain? Same games, budget and seeds.}

\paragraph{Interaction efficiency and compute.} \todo{Learning curves
(Figure~\ref{fig:curves}); DQN vs PPO actions to first level; GPU memory and
time per run.}

\paragraph{Generalisation.} \todo{Dev vs held-out gap. Be explicit: the design
was tuned on dev games only.}

\paragraph{Failure cases.} \todo{Two replays: one success, one characteristic
failure (e.g. agent stuck clicking empty cells, or looping between two
frames). Saved with \texttt{record: true}.}

\paragraph{Limitations.} \todo{Coarse click grid; three seeds; small number of
games; per-game online learning cannot reuse knowledge across games.}
